\documentclass[11pt,a4paper]{amsart}

\usepackage[T1]{fontenc}

\usepackage[english]{babel}

\usepackage{amsmath, amssymb, amsthm}
\usepackage{hyperref}
\usepackage{enumerate}
\usepackage{geometry}
\usepackage{float}
\usepackage{csquotes}

% --- Theorem environments (Ramanujan Journal style) ---
\theoremstyle{plain}
\newtheorem{theorem}{Theorem}[section]
\newtheorem{lemma}[theorem]{Lemma}
\newtheorem{corollary}[theorem]{Corollary}
\newtheorem{proposition}[theorem]{Proposition}
\newtheorem{conjecture}[theorem]{Conjecture}

\theoremstyle{definition}
\newtheorem{definition}[theorem]{Definition}

\theoremstyle{remark}
\newtheorem{remark}[theorem]{Remark}

\usepackage[
    backend=biber,
    style=numeric,
    sorting=none
]{biblatex}
\addbibresource{bibliography.bib}

\geometry{margin=2.8cm}

\title{A Linear-Time Algorithm for the Quadratic Form\\
on the Gap Poset of a Numerical Semigroup,\\
with Generalisation to Multiple Generators}

\author{Artur Flamandzki}
\thanks{Email address: \texttt{flamandzki.artur@gmail.com}}

\date{\today}

\begin{document}

\begin{abstract}
Let $a, b$ be coprime positive integers with $a < b$, and let $G$ be the gap
set of the numerical semigroup $\langle a, b \rangle$, ordered as
$g_0 < g_1 < \cdots < g_{N-1}$ where $N = \tfrac{(a-1)(b-1)}{2}$.
Huang~\cite{Huang2026} introduced a quadratic form
\[
  Q(\mathbf{n}) = \sum_{i,j \in G} K(g_j - g_i)\, n_i\, n_j,
\]
where $K$ is an inclusion--exclusion kernel, and showed that $Q$ recovers
the $\mathrm{dinv}$ statistic from algebraic combinatorics.
The natural evaluation of $Q$ requires $O(N^2)$ arithmetic operations.

We observe that the kernel $K(d)$ is constant on four arithmetic intervals
and vanishes identically for $d < 0$ and $d \geq a+b$.
This arithmetic structure allows the double sum to be condensed into two
sliding-window sums, yielding an $O(N)$ algorithm and reducing complexity
from $O(a^2b^2)$ to $O(ab)$.

We then generalise to numerical semigroups $\langle p_1,\ldots,p_n\rangle$
with $n$ generators. The generalised kernel $K^{(n)}$ is defined by
inclusion--exclusion over all $2^n$ subsets of generators. We show that the
number of active windows $w(n)$ grows strictly slower than $2^n$, due to
cancellations arising from redundant subset sums of mixed parity. The
resulting algorithm runs in $O(w(n) \cdot N)$. While $w(n)$ grows with $n$,
it grows strictly slower than $2^n$: the ratio $w(n)/(2^n - 1)$ decreases in
direct analogy to the Mertens sieve product $\prod(1 - 1/p_i)$
(Theorem~\ref{thm:density}); the gap set size $N$ does not increase when
generators are added. Absolute complexity therefore depends on generator
density --- for small dense generators $N$ collapses, for sparse ones $N$
dominates. The asymptotic behaviour is verified numerically across eight
qualitatively distinct generator families up to $n = 300$.
\end{abstract}

\maketitle

% ==================================================
\section{Introduction}
\label{sec:introduction}
% ==================================================

Let $a, b \in \mathbb{N}$ with $\gcd(a,b) = 1$ and $a < b$.
The \emph{numerical semigroup} generated by $a$ and $b$ is
\[
  \langle a, b \rangle = \{ xa + yb : x, y \in \mathbb{Z}_{\geq 0} \}.
\]
By the Sylvester--Frobenius theorem, the complement
$G = \mathbb{N} \setminus \langle a, b \rangle$
is finite with $|G| = N := \tfrac{(a-1)(b-1)}{2}$.
We write $G = \{g_0, g_1, \ldots, g_{N-1}\}$ with $g_0 < g_1 < \cdots < g_{N-1}$.

In the recent preprint~\cite{Huang2026}, Huang defines a quadratic form
on $\mathbb{R}^G$ in terms of the kernel function
\[
  K(d) := \mathbf{1}_{d \geq 0} - \mathbf{1}_{d \geq a}
         - \mathbf{1}_{d \geq b} + \mathbf{1}_{d \geq a+b}, \quad d \in \mathbb{Z},
\]
and proves that this form recovers the $\mathrm{dinv}$ statistic for rational
Dyck paths. A Lean~4 formalization of this result was carried out by
AxiomProver~\cite{AxiomProver2026}.

Both works address the \emph{correctness} of the identity $Q = \mathrm{dinv}$.
The present note addresses the orthogonal question of computational efficiency
and its generalisation to $n$ generators.

\subsection*{Relation to the Frobenius problem}

The generalisation to $n$ generators does \emph{not} address the
Frobenius problem for $n \geq 3$, which asks for a closed formula
for the largest integer not in $\langle p_1,\ldots,p_n \rangle$
and remains open for $n \geq 3$~\cite{Ramirez2005}.

The present work operates under a fundamentally different question.
The kernel $K^{(n)}$ is defined solely by inclusion--exclusion over
the generator set, and its interval structure follows directly from
this definition, independently of the Frobenius number or the genus
of the semigroup. The sliding-window algorithm takes the sorted gap
set $G$ as input --- it does not compute $G$ from generators.
The computational barrier of the Frobenius problem therefore does
not apply here.

This observation is itself enabled by Huang's formulation: it is
precisely the inclusion--exclusion structure of $K$ that makes the
arithmetic of $K^{(n)}$ transparent and its generalisation direct.
Huang's definition, designed for a geometric purpose, incidentally
exposes a structure that extends well beyond its original scope.

\subsection*{Main results}

\begin{theorem}[Linear evaluation, two generators]
\label{thm:main}
For any $\mathbf{n} \in \mathbb{R}^G$,
\[
  Q(\mathbf{n}) = \sum_{k=0}^{N-1} n_k \bigl(W^+_k - W^-_k\bigr),
\]
where
\begin{align*}
  W^+_k &= \sum_{\substack{j : g_k \leq g_j < g_k + a}} n_j, \\
  W^-_k &= \sum_{\substack{j : g_k + b \leq g_j < g_k + a + b}} n_j.
\end{align*}
Both families can be computed in $O(N)$ total operations via two sliding windows.
\end{theorem}

\begin{theorem}[Generalised sliding-window evaluation]
\label{thm:general}
Let $\mathbf{p} = (p_1,\ldots,p_n)$ with $\gcd(p_1,\ldots,p_n)=1$,
and let $G$ be the gap set of $\langle p_1,\ldots,p_n\rangle$.
Define
\[
  K^{(n)}(d) = \sum_{S \subseteq \mathbf{p}} (-1)^{|S|} \mathbf{1}_{d \geq \sigma(S)},
  \qquad \sigma(S) = \sum_{p \in S} p.
\]
Let $w(n)$ denote the number of intervals on which $K^{(n)}$ is nonzero.
Then $Q(\mathbf{n})$ is computable in time $O(T_{\mathrm{prep}} + w(n) \cdot N)$,
where $T_{\mathrm{prep}} = O(n \cdot \sigma_n)$ via the incremental algorithm
(Theorem~\ref{thm:incremental}).
\end{theorem}

\begin{theorem}[Incremental window computation]
\label{thm:incremental}
The active window set for $n$ generators can be computed incrementally
from the active window set for $n-1$ generators in $O(\sigma_{n-1})$
time, where $\sigma_{n-1} = \sum_{i=1}^{n-1} p_i$.
The total preprocessing cost for all $n$ generators is $O(n \cdot \sigma_n)$,
compared to $O(2^n)$ for direct enumeration. For prime generators
($p_i$ = $i$-th prime), $\sigma_n = \Theta(n^2 \log n)$ by the prime number
theorem, giving $O(n^3 \log n)$.
\end{theorem}

% ==================================================
\section{Arithmetic Structure of the Kernel}
\label{sec:kernel}
% ==================================================

\subsection{Two generators}

\begin{lemma}[Interval decomposition of $K$]
\label{lem:K-intervals}
For all $d \in \mathbb{Z}$:
\[
  K(d) =
  \begin{cases}
    0  & \text{if } d < 0, \\
    +1 & \text{if } 0 \leq d < a, \\
    0  & \text{if } a \leq d < b, \\
    -1 & \text{if } b \leq d < a+b, \\
    0  & \text{if } d \geq a+b.
  \end{cases}
\]
\end{lemma}

\begin{proof}
By direct computation from
$K(d) = \mathbf{1}_{d \geq 0} - \mathbf{1}_{d \geq a} - \mathbf{1}_{d \geq b}
       + \mathbf{1}_{d \geq a+b}$,
using $a < b$:

\medskip
\begin{center}
\begin{tabular}{c|cccc|c}
$d$ & $\mathbf{1}_{d\geq 0}$ & $-\mathbf{1}_{d\geq a}$ &
      $-\mathbf{1}_{d\geq b}$ & $+\mathbf{1}_{d\geq a+b}$ & $K(d)$ \\
\hline
$d < 0$          & 0 & 0 & 0 & 0 & $0$  \\
$0 \leq d < a$   & 1 & 0 & 0 & 0 & $+1$ \\
$a \leq d < b$   & 1 & 1 & 0 & 0 & $0$  \\
$b \leq d < a+b$ & 1 & 1 & 1 & 0 & $-1$ \\
$d \geq a+b$     & 1 & 1 & 1 & 1 & $0$  \\
\end{tabular}
\end{center}
The five cases are exhaustive and mutually exclusive.
\end{proof}

% ==================================================
\section{Reduction of the Double Sum}
\label{sec:reduction}
% ==================================================

\begin{lemma}[Support restriction]
\label{lem:support}
Every term in
\begin{equation}
\label{eq:Q-double}
  Q(\mathbf{n}) = \sum_{k=0}^{N-1} \sum_{\ell=0}^{N-1}
    K(g_\ell - g_k)\, n_k\, n_\ell
\end{equation}
with $g_\ell - g_k < 0$ or $g_\ell - g_k \geq a+b$ contributes zero.
\end{lemma}

\begin{proof}
Immediate from $K(d) = 0$ for $d < 0$ and $d \geq a+b$
(Lemma~\ref{lem:K-intervals}).
\end{proof}

\begin{lemma}[Interval separation]
\label{lem:separation}
For fixed $k$,
\[
  \sum_{\substack{\ell : 0 \leq g_\ell - g_k < a+b}}
    K(g_\ell - g_k)\, n_\ell
  = W^+_k - W^-_k.
\]
\end{lemma}

\begin{proof}
From Lemma~\ref{lem:K-intervals}: $K = +1$ on $[0,a)$ and $K = -1$
on $[b, a+b)$, zero otherwise.
\end{proof}

\begin{proof}[Proof of Theorem~\ref{thm:main}]
Combine Lemmas~\ref{lem:support} and~\ref{lem:separation}.
\end{proof}

% ==================================================
\section{Linear-Time Algorithm}
\label{sec:algorithm}
% ==================================================

\begin{theorem}[Complexity, two generators]
\label{thm:complexity}
$(W^+_k - W^-_k)_{k=0}^{N-1}$, and hence $Q(\mathbf{n})$,
can be computed in $O(N)$ time and $O(N)$ space.
\end{theorem}

\begin{proof}
Maintain four pointers
$\mathrm{lo}^+, \mathrm{hi}^+, \mathrm{lo}^-, \mathrm{hi}^-$
into sorted $G$, initialized to $0$, with running sums $W^+ = 0$
and $W^- = 0$. For each $k = 0, 1, \ldots, N-1$:
\begin{enumerate}[(i)]
\item advance $\mathrm{hi}^+$ while $g_{\mathrm{hi}^+} < g_k + a$,
      adding $n_{\mathrm{hi}^+}$ to $W^+$;
\item retract $\mathrm{lo}^+$ while $g_{\mathrm{lo}^+} < g_k$,
      subtracting $n_{\mathrm{lo}^+}$ from $W^+$;
\item advance $\mathrm{hi}^-$ while $g_{\mathrm{hi}^-} < g_k + a + b$,
      adding $n_{\mathrm{hi}^-}$ to $W^-$;
\item retract $\mathrm{lo}^-$ while $g_{\mathrm{lo}^-} < g_k + b$,
      subtracting $n_{\mathrm{lo}^-}$ from $W^-$;
\item accumulate $n_k \cdot (W^+ - W^-)$.
\end{enumerate}
Each pointer traverses $G$ at most once: total $\leq 4N$ updates.
\end{proof}

\begin{corollary}
\label{cor:complexity}
In terms of the generators, $N = \tfrac{(a-1)(b-1)}{2} = O(ab)$,
so the algorithm runs in $O(ab)$ time versus $O(a^2b^2)$ naive.
\end{corollary}

% ==================================================
\section{Incremental Computation of Active Windows}
\label{sec:incremental}
% ==================================================

The active window set $\mathcal{W}$ for $n$ generators requires
knowing all subset sums of $\{p_1,\ldots,p_n\}$.
Direct enumeration costs $O(2^n)$.
We show this can be reduced to $O(n \cdot \sigma_n)$ where
$\sigma_n = \sum_{i=1}^n p_i$.

\begin{definition}[Parity balance]
\label{def:parity}
For each value $s \in \mathbb{N}$, define the \emph{parity balance}
\[
  \pi_n(s) = \sum_{\substack{S \subseteq \{p_1,\ldots,p_n\} \\ \sigma(S) = s}}
             (-1)^{|S|}.
\]
Then $K^{(n)}(d) = \sum_{s \leq d} \pi_n(s)$ (prefix sum of $\pi_n$).
\end{definition}

\begin{lemma}[Incremental update]
\label{lem:incremental}
Adding generator $p_{n+1}$ updates the parity balance as follows:
\[
  \pi_{n+1}(s) = \pi_n(s) - \pi_n(s - p_{n+1}),
\]
where by convention $\pi_n(s) = 0$ for $s \notin \{0,\ldots,\sigma_n\}$,
and the singleton $\{p_{n+1}\}$ is captured automatically through
$\pi_n(0) = 1$ (the contribution of the empty subset).
\end{lemma}

\begin{proof}
Subsets of $\{p_1,\ldots,p_{n+1}\}$ with sum $s$ are either:
subsets $S \subseteq \{p_1,\ldots,p_n\}$ with $\sigma(S) = s$, contributing
$(-1)^{|S|}$ each --- their sum is $\pi_n(s)$;
or subsets of the form $S \cup \{p_{n+1}\}$ where $S \subseteq \{p_1,\ldots,p_n\}$
with $\sigma(S) = s - p_{n+1}$, contributing $(-1)^{|S|+1} = -(-1)^{|S|}$ each ---
their sum is $-\pi_n(s - p_{n+1})$.
The singleton $\{p_{n+1}\}$ corresponds to $S = \emptyset$ in the second class
at $s = p_{n+1}$, contributing $-\pi_n(0) = -1$, as expected.
\end{proof}

\begin{proof}[Proof of Theorem~\ref{thm:incremental}]
Each update step processes at most $|\mathrm{supp}(\pi_n)| \leq \sigma_n$
values. Summing over $n$ steps gives total cost $\sum_{k=1}^n \sigma_k
= O(n \cdot \sigma_n) = O(n^3 \log n)$ for prime generators.
\end{proof}

\begin{remark}[Benchmark]
The incremental algorithm computes active windows for $n = 2, \ldots, 30$
without the exponential cost of direct subset enumeration.
As a cross-check, direct $O(2^n)$ enumeration was run independently
(Section~\ref{sec:experimental}) and reached $n = 28$, confirming
agreement with the incremental results at $n = 27$ and $n = 28$
on both experimental platforms.
\end{remark}

\begin{remark}
The original implementation of $Q$ in~\cite{Huang2026} and its
Lean~4 formalization~\cite{AxiomProver2026} evaluate the double
sum~\eqref{eq:Q-double} directly, requiring $N^2$ operations per
input vector; the sliding-window algorithm requires $4N$. The ratio
\[
  \frac{N^2}{4N} \;=\; \frac{N}{4} \;=\; \frac{(a-1)(b-1)}{8}
\]
grows linearly with $N$ and is therefore unbounded.
For $(a,b)=(97,101)$ with $N=4801$ this gives $23\,049\,601$ versus
$19\,204$ operations, a ratio of approximately $1\,200$;
systematic numerical studies on the original formulation become
prohibitive at this scale.

Regarding numerical precision: the observed discrepancy of $1.3 \times 10^{-9}$
at $N = 4801$ is not algorithmic disagreement but floating-point accumulation.
The naive $O(N^2)$ sum accumulates rounding errors at rate $\sim N^2 \varepsilon \lvert Q \rvert$,
while the sliding-window algorithm accumulates at rate $\sim N \varepsilon \lvert Q \rvert$,
where $\varepsilon \approx 2.2 \times 10^{-16}$ is machine epsilon.
Both algorithms agree with each other to within $2 \times 10^{-10}$ at $N = 4801$
(Table~\ref{tab:two-gen}; prefix and linear columns).
The theoretical threshold at which naive accumulation reaches $\lvert \text{error} \rvert = 1$
is $N \approx 165\,000$, corresponding to generators $(a,b) \approx (406, 408)$;
for the sliding-window algorithm this threshold is $N \approx 10^8$.
All results in Table~\ref{tab:two-gen} are within full double-precision accuracy.
\end{remark}

% ==================================================
\section{Generalisation to \texorpdfstring{$n$}{n} Generators}
\label{sec:generalisation}
% ==================================================

\subsection{Generalised kernel and active windows}

\begin{definition}[Generalised kernel]
\label{def:K-general}
Let $\mathbf{p} = (p_1, \ldots, p_n)$ with $\gcd(p_1,\ldots,p_n) = 1$.
\[
  K^{(n)}(d) = \sum_{S \subseteq \{p_1,\ldots,p_n\}} (-1)^{|S|}\,
               \mathbf{1}_{d \geq \sigma(S)}, \qquad
  \sigma(\emptyset) = 0.
\]
\end{definition}

\begin{lemma}[Vanishing outside $[0, \sum p_i)$]
\label{lem:K-vanish}
$K^{(n)}(d) = 0$ for $d < 0$ and for $d \geq \sum_{i=1}^n p_i$.
\end{lemma}

\begin{proof}
For $d < 0$: every indicator is zero.
For $d \geq \sum p_i$: all indicators equal~$1$, and
$\sum_{S \subseteq \mathbf{p}} (-1)^{|S|} = (1-1)^n = 0$
by the binomial theorem.
\end{proof}

\begin{definition}[Active window]
\label{def:active}
An interval $[s_k, s_{k+1})$ between consecutive breakpoints of $K^{(n)}$
is \emph{active} if $K^{(n)}$ is nonzero on it.
Let $w(n) = w(\mathbf{p})$ denote the number of active windows.
\end{definition}

\subsection{Redundant subset sums and parity cancellation}

\begin{definition}[Redundant subset sum]
\label{def:redundant}
A value $s \in \mathbb{N}$ is a \emph{redundant subset sum} if
$s = \sigma(S_1) = \sigma(S_2)$ for two distinct subsets $S_1 \neq S_2$.
It is \emph{cancelling} if $|S_1| \not\equiv |S_2| \pmod{2}$.
\end{definition}

\begin{lemma}[Cancellation mechanism]
\label{lem:cancel}
Let $s$ be a cancelling redundant subset sum. Then subsets of opposite
parity contribute opposite signs to $K^{(n)}(s)$, reducing
$|K^{(n)}(s)|$ relative to the non-redundant case and potentially
driving $K^{(n)}(s) = 0$, thereby eliminating a breakpoint from
the active set.
\end{lemma}

\begin{proof}
Direct from Definition~\ref{def:K-general}: each subset $S$ contributes
$(-1)^{|S|}$ at every point $d \geq \sigma(S)$. Two subsets with equal
sum and opposite parities contribute $+1$ and $-1$ at the same breakpoint.
\end{proof}

\begin{lemma}[Linear upper bound on $w(n)$]
\label{lem:w-le-sigma}
For any generator sequence $\mathbf{p} = (p_1,\ldots,p_n)$ with
$\gcd(p_1,\ldots,p_n) = 1$,
\[
  w(n) \;\leq\; \sigma_n \;=\; \sum_{i=1}^n p_i.
\]
\end{lemma}

\begin{proof}
Every breakpoint of $K^{(n)}$ is a subset sum $\sigma(S)$ for some
$S \subseteq \mathbf{p}$, and all such sums lie in $[0, \sigma_n] \cap \mathbb{Z}$,
which contains $\sigma_n + 1$ integers. The active windows are intervals
$[s_k, s_{k+1})$ between consecutive breakpoints with $K^{(n)}$ nonzero;
their number is at most the number of breakpoints minus one, which is at
most $\sigma_n$.
\end{proof}

\begin{corollary}[Asymptotic strict inequality]
\label{cor:subexp-asymp}
For any sequence with $M(n) = \prod (1 - 1/p_i) \to 0$, the ratio
$w(n)/(2^n - 1)$ tends to $0$ asymptotically faster than $M(n)$. In
particular $w(n) < 2^n - 1$ for all sufficiently large $n$.
\end{corollary}

\begin{proof}
Immediate from Theorem~\ref{thm:density}.
\end{proof}

\begin{remark}[Empirical strict inequality at finite $n$]
\label{rmk:subexp-empirical}
Numerical data (Table~\ref{tab:mertens}; eight sequence families to $n=300$
in \texttt{multi\_sequence\_results.txt}) confirm $w(n) < 2^n - 1$ and
$w(n)/(2^n-1)$ monotonically decreasing for every $n \geq 3$ tested,
with the redundancy rate (Lemma~\ref{lem:cancel}) growing from $0\%$
at $n=2$ to $100\%$ from $n = 21$ onward when generators are the first $n$
primes. A combinatorial proof of strict inequality at every finite $n \geq 3$
for arbitrary coprime sequences is left as an open problem.
\end{remark}

\subsection{Algorithm and complexity}

\begin{proof}[Proof of Theorem~\ref{thm:general}]
Compute the active window set $\mathcal{W}$ via the incremental algorithm
of Theorem~\ref{thm:incremental} in time $T_{\mathrm{prep}} = O(n \cdot \sigma_n)$.
For each $(lo, hi, c) \in \mathcal{W}$, run one sliding-window pass
over sorted $G$ in $O(N)$ time. Total cost is therefore
$O(T_{\mathrm{prep}} + w(n) \cdot N)$.
\end{proof}

\begin{remark}[Empirical cost reduction]
\label{rmk:monotone}
Adding generator $p_{n+1}$ produces a new gap set $G'$ with $|G'| \leq |G|$
(adding generators can only enlarge the semigroup), and increases the
window count by $\Delta w(n+1) := w(n+1) - w(n)$. Empirically (see
Section~\ref{sec:open}, item~(vii)) $\Delta w(n) \approx p_n$ for large $n$,
while $w(n) \approx \sigma_n$, so the relative increment $\Delta w(n)/w(n)$
behaves as $1/n$. Both factors of the cost $w(n) \cdot N$ thus tend to
diminish per added generator, in informal analogy to the prime sieve;
a formal asymptotic statement requires Open Problem~(vii).
\end{remark}

% ==================================================
\section{Numerical Verification}
\label{sec:verification}
% ==================================================

\subsection{Experimental setup}
\label{sec:experimental}

Numerical experiments were carried out in Python~3.14 (with SymPy~1.14
and NumPy) on two independent platforms: a workstation with an
AMD~Ryzen~9~5900X processor (12 cores, 24 threads) and 64\,GB of memory;
and a laptop with an Intel Core~i7-13620H processor (6P+4E cores, 16 threads)
and 16\,GB of memory. Both machines run Windows~11~Pro with identical
Python and library versions. Floating-point computation uses IEEE\,754
double precision (machine epsilon $\varepsilon \approx 2.2 \times 10^{-16}$).
Reference scripts: \texttt{verify\_linear\_Q.py}, \texttt{verify\_v2.py},
\texttt{incremental\_windows.py}.

Throughout this section, an \emph{operation} denotes one execution of the
inner-loop body of the corresponding procedure: in the naive $O(N^2)$
algorithm, one term $K(g_\ell - g_k)\,n_k\,n_\ell$ of the double
sum~\eqref{eq:Q-double}; in the sliding-window $O(N)$ algorithm, one
pointer advance with the associated update of the running window sum.
Operation counts are deterministic functions of the input pair $(a,b)$
alone, not wall-clock measurements.

Computations using direct $O(2^n)$ subset enumeration --- employed as a
cross-check against the incremental algorithm at $n = 27$ and $n = 28$ ---
use multi-threaded execution; all other reported quantities are produced
single-threaded and are bit-identical across both platforms.

\subsection{Two generators}

Three implementations --- naive $O(N^2)$, prefix-sum $O(N\log N)$,
and sliding-window $O(N)$ --- were verified on 450 coprime pairs
$(a,b)$ with $b \leq 40$, using 10 random cone vectors each.
All results agreed to tolerance $10^{-9}$ (tolerance $10^{-8}$ for
$N > 500$, accommodating floating-point accumulation).

\begin{table}[H]
\centering
\begin{tabular}{rrrrrl}
\hline
$a$ & $b$ & $N$ & $O(N^2)$ ops & $O(N)$ ops & max error \\
\hline
  3 &   5 &     4 &          16 &        16 & $5.3\times10^{-15}$ \\
  5 &   7 &    12 &         144 &        48 & $1.8\times10^{-14}$ \\
  7 &  11 &    31 &         961 &       124 & $1.1\times10^{-13}$ \\
 11 &  13 &    60 &        3600 &       240 & $2.8\times10^{-13}$ \\
 23 &  29 &   314 &       98596 &      1256 & $1.2\times10^{-11}$ \\
 37 &  41 &   721 &      519841 &      2884 & $4.3\times10^{-11}$ \\
 97 & 101 &  4801 &   23049601 &     19204 & $1.3\times10^{-9}$  \\
\hline
\multicolumn{6}{l}{All 450 pairs: \textbf{PASS}.
  Pairs with $N>500$ verified against $O(N\log N)$ reference.}
\end{tabular}
\caption{Two-generator verification.}
\label{tab:two-gen}
\end{table}

\begin{remark}
The naive $O(N^2)$ implementation directly transcribes the definition
of~\cite{Huang2026} and serves as ground truth for small $N$.
For $(a,b) = (97,101)$ its runtime becomes prohibitive, illustrating
that the original formulation is too costly to benchmark its own
optimisation at scale. The $O(N\log N)$ prefix-sum algorithm, itself
verified against naive for small $N$, closes the verification chain
transitively for large $N$.
\end{remark}

\subsection{Multiple generators}

\begin{table}[H]
\centering
\begin{tabular}{lrrrl}
\hline
Generators & $N$ & Active windows & max error & Status \\
\hline
$\{2,3\}$     & 1 & 2 & $0$                   & OK \\
$\{3,5\}$     & 4 & 2 & $8.9\times10^{-16}$   & OK \\
$\{2,3,5\}$   & 1 & 4 & $0$                   & OK \\
$\{3,5,7\}$   & 3 & 5 & $1.8\times10^{-15}$   & OK \\
$\{2,3,5,7\}$ & 1 & 8 & $0$                   & OK \\
\hline
\end{tabular}
\caption{Multi-generator verification against naive $O(N^2)$.}
\label{tab:multi-gen}
\end{table}

\subsection{Mertens comparison}

\begin{table}[H]
\centering
\small
\begin{tabular}{r|rrrrr}
\hline
$n$ & $2^n-1$ & $w(n)$ & $M(n)$ & Redund.\% & Gap \\
\hline
 2 &          3 &    2 & 0.3333 &  0.0\% & $-0.3333$ \\
 3 &          7 &    4 & 0.2667 & 14.3\% & $-0.3048$ \\
 4 &         15 &    8 & 0.2286 & 26.7\% & $-0.3048$ \\
 5 &         31 &   14 & 0.2078 & 29.0\% & $-0.2438$ \\
 6 &         63 &   20 & 0.1918 & 44.4\% & $-0.1257$ \\
 7 &        127 &   30 & 0.1805 & 59.1\% & $-0.0557$ \\
 8 &        255 &   38 & 0.1710 & 72.2\% & $+0.0220$ \\
 9 &        511 &   48 & 0.1636 & 81.6\% & $+0.0697$ \\
10 &       1023 &   60 & 0.1579 & 88.0\% & $+0.0993$ \\
11 &       2047 &   84 & 0.1529 & 92.5\% & $+0.1118$ \\
12 &       4095 &  116 & 0.1487 & 95.3\% & $+0.1204$ \\
13 &       8191 &  146 & 0.1451 & 97.2\% & $+0.1273$ \\
14 &      16383 &  198 & 0.1417 & 98.3\% & $+0.1296$ \\
15 &      32767 &  236 & 0.1387 & 99.0\% & $+0.1315$ \\
16 &      65535 &  288 & 0.1361 & 99.4\% & $\mathbf{+0.1317}$ \\
17 &     131071 &  352 & 0.1338 & 99.7\% & $+0.1311$ \\
18 &     262143 &  412 & 0.1316 & 99.8\% & $+0.1300$ \\
19 &     524287 &  492 & 0.1296 & 99.9\% & $+0.1287$ \\
20 &    1048575 &  564 & 0.1278 & 99.9\% & $+0.1273$ \\
21 &    2097151 &  646 & 0.1260 &100.0\% & $+0.1257$ \\
22 &    4194303 &  720 & 0.1245 &100.0\% & $+0.1243$ \\
23 &    8388607 &  808 & 0.1230 &100.0\% & $+0.1229$ \\
24 &   16777215 &  892 & 0.1216 &100.0\% & $+0.1215$ \\
25 &   33554431 &  996 & 0.1203 &100.0\% & $+0.1203$ \\
26 &   67108863 & 1096 & 0.1191 &100.0\% & $+0.1191$ \\
27 &  134217727 & 1200 & 0.1180 &100.0\% & $+0.1180$ \\
28 &  268435455 & 1306 & 0.1169 &100.0\% & $+0.1169$ \\
29 &  536870911 & 1416 & 0.1158 &100.0\% & $+0.1158$ \\
30 & 1073741823 & 1528 & 0.1148 &100.0\% & $+0.1148$ \\
\hline
\end{tabular}
\caption{Active window count $w(n)$, Mertens product $M(n)$,
redundancy rate, and $\mathrm{Gap}(n) = M(n) - w(n)/(2^n-1)$,
generators $= p_1,\ldots,p_n$.
Three observed milestones: $n=8$ (smallest $n$ with $\mathrm{Gap}(n) > 0$),
$n=21$ (smallest $n$ with redundancy rate $= 100\%$),
$n=25$ (smallest $n$ with $\mathrm{Gap}(n) = M(n)$ to four decimal places).
Maximum observed Gap at $n=16$ (bold). For $n \geq 25$, $\mathrm{Gap}(n) = M(n)$
within the floating-point precision of our metric.
Data for $n \leq 18$: direct enumeration.
$n = 19,\ldots,30$: incremental algorithm (Section~\ref{sec:incremental}),
verified independently at $n=27,28$ by direct $O(2^n)$ enumeration
(Section~\ref{sec:experimental}).}
\label{tab:mertens}
\end{table}

% ==================================================
\section{Structural Interpretation}
\label{sec:interpretation}
% ==================================================

\subsection{Universality across generator sequences}

The behaviour observed for the first $n$ primes extends to arbitrary
generator sequences. Table~\ref{tab:sequences} summarises results for
eight qualitatively different sequences at $n = 25$.

\begin{table}[H]
\centering
\small
\begin{tabular}{lrrrr}
\hline
Sequence & First gen. & Gap$>0$ at $n$ & Gap$/M$ at $n=25$ & Gap$/M$ at $n=300$ \\
\hline
\multicolumn{5}{l}{\textit{Prime sequences}} \\
First $n$ primes (baseline)        &    2 & 8 & 0.9998 & 1.0000 \\
Primes from $101$                  &  101 & 2 & 0.9999 & 1.0000 \\
Primes from $1009$                 & 1009 & 2 & 0.9997 & 1.0000 \\
Every 2nd prime ($p_2,p_4,\ldots$) &    3 & 8 & 0.9998 & 1.0000 \\
Every 3rd prime ($p_3,p_6,\ldots$) &    5 & 2 & 0.9998 & 1.0000 \\
Lower twin primes                  &    3 & 9 & 0.9996 & 1.0000 \\
Primes $\equiv 3\pmod{4}$          &    3 & 8 & 0.9998 & 1.0000 \\
Primes $\equiv 1\pmod{4}$          &    5 & 2 & 0.9998 & 1.0000 \\
\multicolumn{5}{l}{\textit{Random sequences (GPT density test)}} \\
Random, same density (seed 42)     &  --- & 8 & 0.9999 & --- \\
Random, same density (seed 137)    &  --- & 8 & 0.9999 & --- \\
Random, no density match (seed 42) &  --- & 2 & 0.9999 & --- \\
\hline
\end{tabular}
\caption{Theorem~\ref{thm:density} ($\mathrm{Gap}/M \to 1$)
holds for all tested sequences.
At $n=300$, Gap$/M(n) = 1.0000$ (six decimal places) for all prime sequences.
Random odd integers with the same density as primes give identical
Gap$/M$ and crossover behaviour to the prime baseline ---
confirming the effect is a property of generator \emph{density},
not arithmetic structure.
Random integers without density matching behave like large-prime sequences
(immediate crossover at $n=2$).}
\label{tab:sequences}
\end{table}

\begin{remark}[Density determines crossover]
The crossover point (Gap $> 0$) and the rate of Gap$/M \to 1$
depend on the \emph{density} of generators, not their arithmetic nature.
Small generators (density: many generators per unit interval) delay
the crossover; large generators (sparse density) produce immediate
crossover at $n=2$.
Random odd integers with the same density as the prime sequence
reproduce the prime baseline exactly ($n^* = 8$, Gap$/M = 0.9999$ at $n=30$).
This indicates that the mechanism is combinatorial rather than
number-theoretic: the prime sequence is one example of a sequence
with density $\sim 1/\log n$, and the empirical behaviour depends only
on this density.
\end{remark}

The efficiency gain has a transparent arithmetic explanation.
The kernel $K^{(n)}$ is the inclusion--exclusion M\"{o}bius function
over the generator set. Its sparsity reflects two independent effects.

\begin{enumerate}[(i)]
\item \emph{Mertens-type sparsity.} Each generator $p_i$ eliminates
a fraction $\sim 1/p_i$ of the gap set, reducing $N$.
This accounts for the factor $M(n) = \prod(1 - 1/p_i)$.

\item \emph{Redundancy cancellation.} Distinct subsets with equal sums
and opposite parities annihilate breakpoints in $K^{(n)}$,
reducing $w(n)$ below what Mertens predicts. This effect grows
with $n$ (redundancy rate $0\% \to 99.8\%$ for $n = 2 \to 18$)
and has no analogue in the prime sieve, where primes are by
construction non-redundant.
\end{enumerate}

The combination gives total complexity $O(w(n) \cdot N)$ decreasing
faster than either factor alone.

\subsection{Activation and deactivation of window positions}

Each new generator $p_{n+1}$ modifies $K^{(n+1)}$ relative to $K^{(n)}$
by adding the contribution of every subset containing $p_{n+1}$.

\begin{lemma}[Activation mechanism]
\label{lem:activation}
Let $\mathbf{p}' = (p_1,\ldots,p_n,p_{n+1})$.
For any $d \in \mathbb{Z}$,
\[
  K^{(n+1)}(d) - K^{(n)}(d)
  = \sum_{\substack{S \subseteq \{p_1,\ldots,p_n\} \\ \sigma(S) + p_{n+1} \leq d}}
    (-1)^{|S|+1}.
\]
In particular, a position $d$ that is inactive under $K^{(n)}$ becomes
active under $K^{(n+1)}$ if and only if the above sum is nonzero,
i.e.\ the new generator introduces an imbalance between even- and
odd-sized subsets whose augmented sum falls at or below $d$.
\end{lemma}

\begin{proof}
Direct from Definition~\ref{def:K-general}: subsets not containing
$p_{n+1}$ contribute identically to $K^{(n+1)}$ and $K^{(n)}$.
The difference is the contribution of subsets $S \cup \{p_{n+1}\}$,
which equals $(-1)^{|S|+1} \cdot \mathbf{1}_{d \geq \sigma(S) + p_{n+1}}$.
\end{proof}

\begin{remark}
The lemma shows that positions deactivated at step $n$ can be reactivated
at step $n+1$ if the new generator restores parity imbalance at that point.
Empirically (Table~\ref{tab:positions}), the first few generators are smallest
primes and their pairwise sums are small, so reactivation is frequent for
small positions and rare for large ones.
Moreover, $K^{(n)}(d)$ is not restricted to $\{-1, 0, +1\}$:
values $\pm 2$ and $\pm 3$ are observed at $n = 12$ (e.g.\ positions
$47$ and $149$), indicating that multiple imbalanced subsets can accumulate
at the same breakpoint. No nontrivial bound on $|K^{(n)}(d)|$ is currently
known; this is a natural open question.
\end{remark}

\begin{table}[H]
\centering
\small
\begin{tabular}{r|rrrrrrrrrrr}
\hline
pos & $n{=}2$ & $n{=}3$ & $n{=}4$ & $n{=}5$ & $n{=}6$ &
      $n{=}7$ & $n{=}8$ & $n{=}9$ & $n{=}10$ & $n{=}11$ & $n{=}12$ \\
\hline
 0 &  1 &  1 &  1 &  1 &  1 &  1 &  1 &  1 &  1 &  1 &  1 \\
 3 & -1 & -1 & -1 & -1 & -1 & -1 & -1 & -1 & -1 & -1 & -1 \\
 5 &  · & -1 & -1 & -1 & -1 & -1 & -1 & -1 & -1 & -1 & -1 \\
 7 &  · &  · & -1 & -1 & -1 & -1 & -1 & -1 & -1 & -1 & -1 \\
 8 &  · &  1 &  · &  · &  · &  · &  · &  · &  · &  · &  · \\
 9 &  · &  · &  1 &  1 &  1 &  1 &  1 &  1 &  1 &  1 &  1 \\
10 &  · &  · &  1 &  1 &  1 &  1 &  1 &  1 &  1 &  1 &  1 \\
16 &  · &  · &  · &  · &  1 &  1 &  1 &  1 &  1 &  1 &  1 \\
17 &  · &  · &  · &  1 &  2 &  1 &  1 &  1 &  1 &  1 &  1 \\
25 &  · &  · &  · &  · &  · &  · &  1 &  1 &  1 &  1 &  1 \\
47 &  · &  · &  · &  · &  · &  1 &  1 &  · & -1 & -2 & -3 \\
\hline
\end{tabular}
\caption{Selected active window positions and their $K^{(n)}$ values
for $n = 2,\ldots,12$ (generators $= p_1,\ldots,p_n$).
A dot denotes $K^{(n)}(d) = 0$ (inactive).
Position $8$ activates at $n=3$ and deactivates at $n=4$.
Position $17$ reaches $K=2$ at $n=6$ (amplification).
Position $47$ reaches $K=-3$ at $n=12$.}
\label{tab:positions}
\end{table}

\subsection{Density asymptotic theorem}
\label{sec:density-theorem}

Table~\ref{tab:mertens} reveals that $\mathrm{Gap}(n)$ is not monotone.
It increases from $n=8$ to a numerical maximum near $n=16$
($\approx 0.1317$), then decreases slowly. Empirically, by $n=28$ the
ratio $\mathrm{Gap}(n)/M(n)$ reaches $1.0000$ to four decimal places
and remains there throughout the tested range $n=2,\ldots,300$ across
all eight sequence families of Section~\ref{sec:verification}. The
following theorem explains and formalises this behaviour.

\begin{theorem}[Density asymptotic]
\label{thm:density}
Let $\mathbf{p} = (p_1, p_2, \ldots)$ be a sequence of distinct integers
greater than $1$ with $\gcd(p_1,\ldots,p_n)=1$ and
$M(n) = \prod_{i=1}^n (1-1/p_i) \to 0$ as $n \to \infty$. Assume in
addition the subexponential growth condition $\log p_n = o(n)$
(equivalently $p_n = e^{o(n)}$). Then
\[
  \frac{w(n)}{2^n - 1} = o(M(n)) \quad \text{as } n \to \infty,
\]
equivalently $\mathrm{Gap}(n)/M(n) \to 1$.

The growth condition $\log p_n = o(n)$ is satisfied by every generator
family considered in Section~\ref{sec:verification}: the first $n$ primes
($p_n \sim n \log n$), shifted primes, twin primes, arithmetic progressions
of primes, and random odd integers with prime-like density.
\end{theorem}

\begin{proof}
We bound $w(n)/((2^n-1) M(n))$ and show it tends to zero.

\textbf{Step 1.} $w(n) \leq \sigma_n := \sum_{k=1}^n p_k$ by
Lemma~\ref{lem:w-le-sigma}. The bound is verified empirically to
$n = 30$ (Table~\ref{tab:mertens}; see also \texttt{verify\_v2.py}).

\textbf{Step 2: bound on $H_n = \sum_{k=1}^n 1/p_k$.}
Since $p_1,\ldots,p_n$ are distinct integers $\geq 2$, the harmonic-type
sum $H_n$ is maximised when $\{p_k\}$ is the set $\{2, 3, \ldots, n+1\}$
of the $n$ smallest such integers. Hence
\[
  H_n \;\leq\; \sum_{j=2}^{n+1} \frac{1}{j} \;\leq\; \log(n+1).
\]
By $\log(1-x) \geq -2x$ for $x \in [0, 1/2]$ and $1/p_k \leq 1/2$,
\[
  M(n) \;\geq\; e^{-2H_n} \;\geq\; (n+1)^{-2}.
\]

\textbf{Step 3: combining the bounds.}
Using $\sigma_n \leq n \cdot p_n$ (where $p_n = \max_k p_k$) and
$2^n - 1 \geq 2^{n-1}$,
\[
  \frac{w(n)}{(2^n-1)\,M(n)} \;\leq\; \frac{2\,\sigma_n}{2^n \cdot M(n)}
  \;\leq\; \frac{2\,n\,p_n\,(n+1)^2}{2^n}
  \;=\; 2\,n\,(n+1)^2 \cdot \frac{p_n}{2^n}.
\]

\textbf{Step 4: subexponential growth implies convergence.}
Under the hypothesis $\log p_n = o(n)$,
\[
  \log\!\left(\frac{p_n}{2^n}\right) \;=\; \log p_n - n \log 2
  \;=\; o(n) - n \log 2 \;\to\; -\infty,
\]
so $p_n / 2^n \to 0$ exponentially fast. Multiplying by the polynomial
factor $2 n (n+1)^2$ leaves the right-hand side tending to zero.

Therefore $w(n)/((2^n-1)\,M(n)) \to 0$, equivalently
$w(n)/(2^n-1) = o(M(n))$, and
$\mathrm{Gap}(n)/M(n) = 1 - w(n)/((2^n-1)M(n)) \to 1$.
\end{proof}

\begin{remark}[Density interpretation]
\label{rmk:density-interpretation}
The proof shows that the crossover point $n^*$ (where Gap $> 0$) and
the rate of convergence Gap$/M \to 1$ depend on the \emph{harmonic
density} $H_n = \sum 1/p_k$ of the generators:
large generators contribute small $1/p_k$, so $H_n$ grows slowly,
redundancy saturates quickly, and $n^*$ is small.
Small generators contribute large $1/p_k$, requiring more steps
before the exponential decay of $2^{-n}$ dominates.
This is the formal expression of the threshold between
\emph{density-driven} and \emph{count-driven} regimes.
\end{remark}

The behaviour of $\mathrm{Gap}(n)$ at finite $n$ is consistent with
diminishing increments: each new prime $p_{n+1}$ contributes a
fraction $\sim 1/p_{n+1}$ of additional redundancy relative to the
existing structure, so increments to $\mathrm{Gap}(n)$ decrease at a
rate governed by the growth of $p_n$. By the prime number theorem
$p_n \sim n \log n$, increments decay faster than $1/(n \log n)$,
which is summable --- consistent with the empirical convergence
visible in Table~\ref{tab:mertens}. Whether $\mathrm{Gap}(n)$
converges to a finite limit, reaches a true maximum at some finite
$n$, or merely approaches its limit from above after a turning point,
is not determined by the asymptotic statement and remains an open
question (Section~\ref{sec:open}, item~(iv)).

% ==================================================
\section{Conclusion}
\label{sec:conclusion}
% ==================================================

The present work makes four contributions.

\textbf{(i) Linear-time evaluation of $Q$.}
The quadratic form $Q$ of Huang~\cite{Huang2026} admits an $O(N)$
evaluation algorithm, reducing the naive $O(N^2)$ cost by a factor
of $N$ through direct arithmetic analysis of the interval structure
of the kernel $K$.

\textbf{(ii) Generalisation to $n$ generators.}
The kernel $K^{(n)}$ defined by inclusion--exclusion over $n$ generators
has an interval structure that extends the two-generator case directly.
This generalisation does not require solving the Frobenius problem for
$n \geq 3$: the algorithm takes the gap set $G$ as input and operates
on its elements, independently of how $G$ was computed.
It is precisely Huang's inclusion--exclusion formulation of $K$ that
makes this generalisation transparent --- a structural property of his
definition that was not the focus of the original work.

\textbf{(iii) Incremental preprocessing.}
Active windows can be maintained incrementally as generators are added,
reducing preprocessing from $O(2^n)$ to $O(n \cdot \sigma_n)$.
This makes the algorithm practical for $n$ up to several hundred.

\textbf{(iv) Empirical exploration and asymptotic theorem.}
Numerical data to $n = 300$ (Section~\ref{sec:verification}) for
eight generator families (primes, shifted primes, twin primes,
arithmetic progressions of primes, and random odd integers with
prime-like density) exhibit three observed milestones ($n=8$,
$n=21$, $n=25$) and a parity-separated structure in the increments
$\Delta w(n)$ for $n \geq 19$. Theorem~\ref{thm:density} proves that
$\mathrm{Gap}(n)/M(n) \to 1$ under the subexponential growth condition
$\log p_n = o(n)$, which is satisfied by all tested families. Across
all tested sequences only the empirical crossover point $n^*$ varies;
the asymptotic behaviour is identical.

% ==================================================
\section{Open Problems}
\label{sec:open}
% ==================================================

The following questions arise naturally from the numerical observations
and the activation mechanism of Lemma~\ref{lem:activation}.
They are stated precisely to invite further investigation.

\begin{enumerate}[(i)]

\item \textbf{Unboundedness of $K^{(n)}$.}
Empirically, $|K^{(n)}(d)|$ takes values $2$ and $3$ for $n \leq 12$.
Is $\sup_{n,d} |K^{(n)}(d)|$ finite or infinite?
More specifically: for fixed $d$, does $K^{(n)}(d)$ stabilise as
$n \to \infty$, or does it grow without bound?

\item \textbf{Stable kernel of positions.}
Empirically, a growing set of positions $d$ appear to stabilise ---
their $K^{(n)}(d)$ value remains constant for all sufficiently large $n$.
What is the mechanism that causes a position to stabilise?
Is stabilisation permanent once it occurs, or can a position
destabilise again for larger $n$?

\item \textbf{Crossover threshold $n^*$.}
Theorem~\ref{thm:density} (Section~\ref{sec:density-theorem}) proves
$\mathrm{Gap}(n)/M(n) \to 1$ but does not characterise the crossover
point $n^* = \min\{n : \mathrm{Gap}(n) > 0\}$. Empirically $n^*$ depends
on the harmonic density $H_n = \sum 1/p_k$: sequences with small
generators (large $1/p_k$) have large $n^*$; sequences with large
generators have $n^* = 2$. Is there a closed-form expression for $n^*$
in terms of $H_n$ or the growth rate of $p_k$?

\item \textbf{Maximum of $\mathrm{Gap}(n)$.}
Empirically, $\mathrm{Gap}(n)$ for the first $n$ primes attains a
numerical maximum near $n=16$ ($\approx 0.1317$) and decreases monotonically
thereafter (Table~\ref{tab:mertens}). Does $\mathrm{Gap}(n)$ attain a
true maximum at some finite $n$? Equivalently: is the convergence
$\mathrm{Gap}(n)/M(n) \to 1$ from below for all sufficiently large $n$,
or does the difference oscillate?

\item \textbf{Bipartite structure of $\Delta w(n)$.}
From $n=19$ onward, the increments $\Delta w(n) = w(n) - w(n-1)$
split into two linearly growing sequences according to the parity of $n$:

\begin{center}
\begin{tabular}{r|rrrrrrrr}
\hline
$n$ & 19 & 20 & 21 & 22 & 23 & 24 & 25 & 26 \\
\hline
$\Delta w$ & 80 & 72 & 82 & 74 & 88 & 84 & 104 & 100 \\
parity & odd & even & odd & even & odd & even & odd & even \\
\hline
\end{tabular}
\end{center}

Odd-$n$ increments: $80, 82, 88, 104, 110, \ldots$ (increasing).
Even-$n$ increments: $72, 74, 84, 100, 106, \ldots$ (increasing).
Both sequences grow approximately linearly, offset by $\approx 8$.
This structure does not appear for $n < 19$ and coincides with
the onset of redundancy saturation.
What is the arithmetic mechanism behind this bipartite split?

\item \textbf{Structure of position renewal.}
Lemma~\ref{lem:activation} identifies when a new generator can
reactivate a previously deactivated position.
Is there a characterisation of which positions are renewed at each
step, in terms of the arithmetic of the generators?

\item \textbf{Asymptotic increment law $\Delta w(n) \sim p_n$.}
Empirical data reveals a striking pattern at large $n$:
the increment $\Delta w(n) = w(n) - w(n-1)$ converges to the
generator value $p_n$ itself.

For the baseline sequence (first $n$ primes) at $n=300$:
$p_{300} = 1987$, $\Delta w(300) = 1986$ (difference: $1$).
For the sequence starting from $101$ at $n=300$:
$p_{300} = 2153$, $\Delta w(300) = 2153$ (difference: $0$).

The pattern $\Delta w(n) \approx p_n$ appears to hold for all tested
sequences once $n$ is sufficiently large (and exactly for sequences
of large generators from the first step).

The proposed mechanism: when $p_n$ is large relative to the existing
subset sums, almost every subset $S \cup \{p_n\}$ produces a unique
new breakpoint $\sigma(S) + p_n$ with no parity collision, contributing
a net new active window. The number of such subsets is approximately
$2^{n-1}$ but most cancel in pairs, leaving a residue of order $p_n$.

Formally: does $\Delta w(n) / p_n \to 1$ as $n \to \infty$
for sequences where $p_n \to \infty$?
If so, this would give $w(n) \sim \sum_{k=1}^n p_k = \sigma_n$,
providing an explicit asymptotic for $w(n)$ in terms of the
generator sum.

\end{enumerate}

\subsection*{Remark on related literature}

The cancellation mechanism identified here appears independently in two
other bodies of work, noted for context rather than as direct antecedents.

Whitney~\cite{Whitney1932} observed that terms in inclusion--exclusion
for graph chromatic polynomials cancel predictably; this was generalised
by Chen and Qian~\cite{ChenQian2018} to ordering-free cancellations.
Those results concern graph-theoretic objects and do not address arithmetic
kernels on numerical semigroups.

In analytic number theory, the exponential growth of inclusion--exclusion
terms in sieve theory and the partial cancellations due to parity constraints
are discussed by Tao~\cite{Tao2007} in the context of the parity problem.
The connection to the present setting is structural: the Mertens-type
sparsity of active windows and the redundancy cancellations we observe
are arithmetically analogous to sieve parity phenomena, though the objects
and proof methods differ.

We arrived at this connection by examining the arithmetic behaviour of
$K^{(n)}$ directly, not by importing sieve-theoretic tools.
The parallel emerged \emph{a posteriori} during a search of related
literature on inclusion--exclusion cancellations. To our knowledge,
neither the interval-sparsity of $K^{(n)}$ nor the parity-cancellation
mechanism for numerical semigroup kernels has been previously described.

Source code for all numerical verifications is available as
\texttt{verify\_linear\_Q.py}, \texttt{appendix\_verify.py},
and \texttt{verify\_v2.py}.

\printbibliography

\end{document}
